\section{Clasificadores supervisados}
\label{sec:met_supervisados}
% Extensión estimada: 3.5--4 páginas, con tres tablas, en cuatro subsecciones.
% Fuentes: src/models/classifiers.py; src/data/transforms.py; src/training/recipes.py, optim_groups.py,
% schedulers.py, losses.py; src/evaluation/tta.py; scripts/train_classifier.py; scripts/toko/submit_sweep.py
% y submit.py (valores efectivos); scripts/sweep_utils.py y configs/sweeps/phase1_modern_sweep_grid.yaml
% (grilla); logs de entrenamiento (parámetros); docs/resul/split_lesion_resul_20260926.md (ganadores, semillas,
% réplica del desbalance); docs/resul/densenet_vit_moderna_resul_20261017.md (DenseNet121 y ViT-B/16 con la
% configuración de su familia, tres semillas); scripts/lesion_split_report.py, classifier_ensemble.py,
% melanoma_detection.py. No se usan como fuente los configs/*.yaml de clasificación, porque train_classifier.py no los lee.
% Revisión del autor del 18/10/2026: solo la configuración final; sin comparar recetas ni describir la búsqueda
% anterior de DenseNet121 y ViT-B/16, cuyas corridas quedan fuera del capítulo.

Los clasificadores supervisados cumplen dos funciones: son la referencia con la que se juzga si un \gls{vlm} clasifica
bien (Sección~\ref{sec:met_lineas_base}) y, dentro del \emph{pipeline}, son los que fijan el diagnóstico.

\subsection{Arquitecturas}
\label{sec:met_supervisados_arquitecturas}

Se comparan cuatro arquitecturas, dos convolucionales, DenseNet121 y ConvNeXt-Tiny (Sección~\ref{sec:marco_cnn}), y
dos basadas en transformers, ViT-B/16 y Swin-Tiny (Secciones~\ref{sec:marco_vit} y~\ref{sec:marco_swin}). En cada
familia hay un modelo anterior y uno más reciente, lo que permite comparar no solo los dos tipos de arquitectura sino
también dos generaciones de cada uno. Además, los cuatro parten de pesos preentrenados en \emph{ImageNet}-1K, y su última
capa se reemplaza por una cabeza de siete salidas precedida por un \emph{dropout} de 0.3
(Tabla~\ref{tab:met_modelos_sup}). Todos se entrenan en una GPU de consumo. Respecto del tamaño, las variantes Tiny, con
unos 27.5 millones de parámetros, tienen un tercio de los de ViT-B/16.

\begin{table}[H]
    \centering
    \caption{Los cuatro clasificadores: familia, preentrenamiento, cabeza, parámetros con cabeza de siete clases y tamaño de lote.}
    \label{tab:met_modelos_sup}
    \small
    \begin{tabularx}{\textwidth}{@{}>{\raggedright\arraybackslash}p{2.1cm}>{\raggedright\arraybackslash}X>{\raggedright\arraybackslash}p{1.9cm}>{\raggedright\arraybackslash}p{2.6cm}>{\raggedleft\arraybackslash}p{1.7cm}>{\raggedleft\arraybackslash}p{0.8cm}@{}}
        \toprule
        \textbf{Modelo} & \textbf{Familia} & \textbf{Pesos iniciales} & \textbf{Cabeza} & \textbf{Parámetros} & \textbf{Lote} \\
        \midrule
        DenseNet121 & Red convolucional, bloques densos & \emph{ImageNet}-1K & \emph{Dropout} 0.3 y lineal de 7 salidas & 6\,961\,031 & 32 \\
        \filasep
        ConvNeXt-Tiny & Red convolucional moderna & \emph{ImageNet}-1K & \emph{Dropout} 0.3 y lineal de 7 salidas & 27\,825\,511 & 32 \\
        \filasep
        ViT-B/16 & Transformer visual, parches de $16 \times 16$ & \emph{ImageNet}-1K & \emph{Dropout} 0.3 y lineal de 7 salidas & 85\,804\,039 & 16 \\
        \filasep
        Swin-Tiny & Transformer jerárquico, con ventanas & \emph{ImageNet}-1K & \emph{Dropout} 0.3 y lineal de 7 salidas & 27\,524\,737 & 16 \\
        \bottomrule
    \end{tabularx}
\end{table}

\subsection{Preprocesamiento y aumento de datos}
\label{sec:met_supervisados_preprocesamiento}

En cuanto a la entrada, las imágenes de \emph{HAM10000} se reducen una sola vez a $224 \times 224$ píxeles, sin
conservar su proporción, y todos los modelos de este capítulo parten de esa versión. Los clasificadores, además, las
reciben normalizadas con la media y el desvío de \emph{ImageNet}. Además, con 6\,447 imágenes de entrenamiento (Sección~\ref{sec:met_particiones}), el riesgo
principal es el sobreajuste (Sección~\ref{sec:marco_regularizacion}), y para reducirlo el entrenamiento aplica aumentos
de datos que cambian la imagen sin cambiar su clase, ya que una misma lesión puede fotografiarse con otro encuadre, otra
orientación u otra iluminación. En concreto, se toma un recorte aleatorio de entre el 70\,\% y el 100\,\% del área, se voltea la imagen
horizontalmente con probabilidad 0.5, se aplican dos transformaciones de RandAugment~\cite{cubuk2020randaugment} con una
magnitud $M$ que se eligió por búsqueda, y se borra un rectángulo al azar con probabilidad
0.25~\cite{zhong2020random}.

En cambio, en validación y en prueba la imagen se redimensiona a $256 \times 256$ y se recorta el centro de
$224 \times 224$. Las probabilidades finales de ambas particiones se obtienen con \gls{tta}
(Sección~\ref{sec:marco_regularizacion}), con la imagen y su reflejo horizontal, mientras que la validación de cada
época, que decide la detención temprana, usa una sola vista. Así, la salida de cada clasificador es una probabilidad por
clase, de la que se calculan las métricas de la Sección~\ref{sec:met_protocolo}.

\subsection{Entrenamiento e hiperparámetros}
\label{sec:met_supervisados_receta}

Una vez fijadas las entradas, cada modelo se ajusta en dos etapas, una forma de transferencia con ajuste fino
(Sección~\ref{sec:marco_transferencia}).
En la primera, con el resto de la red congelado, se entrena solo la cabeza durante una época. En la segunda se
descongela toda la red y cada capa recibe 0.85 veces la tasa de aprendizaje de la capa siguiente
(Ecuación~\ref{eq:marco_llrd}). Esta etapa dura hasta 39 épocas, con tres de calentamiento lineal y decaimiento coseno de la tasa
\cite{goyal2017accurate,loshchilov2017sgdr}. Por otra parte, el optimizador es AdamW~\cite{loshchilov2019decoupled} y la
pérdida es la entropía cruzada con suavizado de etiquetas de 0.1~\cite{szegedy2016rethinking,muller2019does}.
Finalmente, la detención temprana corta el entrenamiento si el F1 ponderado de validación no mejora en cinco épocas, y se
conserva la mejor época (Sección~\ref{sec:met_seleccion}). La Tabla~\ref{tab:met_receta} reúne estos valores.

\begin{table}[H]
    \centering
    \caption{Parámetros de entrenamiento comunes a los cuatro clasificadores.}
    \label{tab:met_receta}
    \small
    \begin{tabularx}{\textwidth}{@{}>{\raggedright\arraybackslash}p{4.2cm}>{\raggedright\arraybackslash}X@{}}
        \toprule
        \textbf{Parámetro} & \textbf{Valor} \\
        \midrule
        Optimizador & AdamW \\
        \filasep
        Épocas & Hasta 40: una de cabeza y hasta 39 de red completa \\
        \filasep
        Decaimiento de la tasa por capa & 0.85 \\
        \filasep
        Planificador & Calentamiento lineal de tres épocas desde 0.1 de la tasa; luego coseno \\
        \filasep
        Pérdida & Entropía cruzada, con suavizado de etiquetas de 0.1 \\
        \filasep
        Detención temprana & Paciencia de cinco épocas sobre el F1 ponderado de validación \\
        \filasep
        Aumento de datos & Recorte aleatorio (70--100\,\% del área), volteo horizontal, RandAugment ($N = 2$) y borrado aleatorio ($p = 0.25$) \\
        \filasep
        Inferencia & \gls{tta} con la imagen y su reflejo horizontal \\
        \filasep
        Precisión numérica & \emph{bfloat16} (bf16) \\
        \filasep
        Semillas & 42 en la búsqueda; 42, 43 y 44 en la configuración elegida de cada modelo \\
        \bottomrule
    \end{tabularx}
\end{table}

Además de estos valores comunes, tres hiperparámetros se eligieron por búsqueda: la tasa de aprendizaje, el decaimiento
de pesos y la magnitud de RandAugment (Tabla~\ref{tab:met_barrido}). Las cinco tasas se cruzaron con los tres decaimientos, y cada magnitud se
asignó a cinco de esas 15 combinaciones. Todas se entrenaron con la semilla 42, con y sin pesos por clase en la
pérdida (Sección~\ref{sec:marco_desbalance}), lo que suma 30 configuraciones por modelo. La
búsqueda se hizo sobre Swin-Tiny y ConvNeXt-Tiny, y en ambos ganó, por F1 macro en validación, una configuración sin
pesos por clase. En cambio, DenseNet121 y ViT-B/16 no tuvieron búsqueda propia, sino que cada uno usa la configuración
ganadora del modelo de su misma familia, la de ConvNeXt-Tiny en el caso de DenseNet121 y la de Swin-Tiny en el de ViT-B/16. Así los cuatro
se comparan con el mismo procedimiento, aunque con una salvedad: si DenseNet121 o ViT-B/16 rinden menos, puede deberse a
que la configuración heredada no es la mejor para ellos. La grilla completa está en el Apéndice~\ref{ap:hiperparametros}.

\begin{table}[H]
    \centering
    \caption{Hiperparámetros explorados en la búsqueda y valores elegidos. DenseNet121 usa los de ConvNeXt-Tiny, y ViT-B/16, los de Swin-Tiny.}
    \label{tab:met_barrido}
    \small
    \begin{tabularx}{\textwidth}{@{}>{\raggedright\arraybackslash}p{3.4cm}>{\raggedright\arraybackslash}X>{\raggedright\arraybackslash}p{2.0cm}>{\raggedright\arraybackslash}p{2.2cm}@{}}
        \toprule
        \textbf{Parámetro} & \textbf{Valores explorados} & \textbf{Swin-Tiny} & \textbf{ConvNeXt-Tiny} \\
        \midrule
        Tasa de aprendizaje & $3.0 \times 10^{-5}$; $7.5 \times 10^{-5}$; $1.5 \times 10^{-4}$; $2.5 \times 10^{-4}$; $5.0 \times 10^{-4}$ & $7.5 \times 10^{-5}$ & $3.0 \times 10^{-5}$ \\
        \filasep
        Decaimiento de pesos & 0.01; 0.05; 0.10 & 0.10 & 0.05 \\
        \filasep
        Magnitud de RandAugment ($M$) & 6; 9; 12, cada valor en cinco de las 15 combinaciones & 6 & 9 \\
        \filasep
        Pesos por clase & Inversos a la frecuencia; ninguno & Ninguno & Ninguno \\
        \bottomrule
    \end{tabularx}
\end{table}

El tratamiento del desbalance merece una comparación aparte, porque los pesos por clase no son la única forma de
abordarlo (Sección~\ref{sec:marco_desbalance}). Para ello, sobre Swin-Tiny, ConvNeXt-Tiny y DenseNet121 se comparan
cuatro pérdidas: la entropía cruzada sin pesos, la ponderada con pesos inversos a la frecuencia de cada clase, la
pérdida focal con $\gamma = 2$~\cite{lin2017focal}, que se combina con esos mismos pesos, y el \emph{balanced
softmax}~\cite{ren2020balanced}, sin pesos.
Para que la pérdida sea lo único que cambia, esta comparación fija el resto de los
hiperparámetros (decaimiento de pesos 0.05, $M = 9$ y una tasa de $10^{-4}$, o de $7.91 \times 10^{-5}$ en DenseNet121)
y usa una sola semilla.

\subsection{Ensamble y detección de melanoma}
\label{sec:met_supervisados_ensamble}

Además de usarse por separado, los dos modelos de la búsqueda, Swin-Tiny y ConvNeXt-Tiny, se combinan por voto suave
(Sección~\ref{sec:marco_senales_confianza}). Solo se combinan esos dos porque son los que eligió la validación. Elegir otra
combinación exigiría consultar la prueba. Por otra parte, para la detección de melanoma, el puntaje es la
probabilidad de melanoma de cada modelo o del ensamble, con el umbral de la Sección~\ref{sec:met_seleccion}.

Estos clasificadores devuelven una probabilidad por clase, pero no generan texto. Para saber si un \gls{vlm}, que sí lo
genera, clasifica tan bien como ellos, se los compara con las mismas métricas (Sección~\ref{sec:met_clasif_vlm}).
